% Self-contained mathematical body; requires amsmath, amssymb, amsthm,
% hyperref; theorem, lemma, proposition, corollary, remark environments.
% Suggested title: The shortest positive average cancelling two frequencies:
% exact classification and a cubic-root resonance law.

\section{The extremal question and the result}

For a compactly supported probability measure $\mu$ on $\mathbb R$ write
\[
 \widehat\mu(\omega)=\int_{\mathbb R}e^{i\omega u}\,d\mu(u),
 \qquad
 \mathcal B(r)=\inf\{\operatorname{diam}(\operatorname{supp}\mu):
       \widehat\mu(1)=\widehat\mu(r)=0\},\quad r>1.
\]
Translations preserve both vanishing conditions.  Thus a measure of support
diameter $2T$ can always be placed in $[-T,T]$.  No symmetry of the admissible
measure is assumed.  The theorem below determines the infimum, all minimizers,
and two scalar equations computing them.  Irrational frequency ratios are
allowed.

\begin{theorem}[Exact two-frequency classification]\label{filt:classification}
Let $r>1$.
\begin{enumerate}
\item If $r=n\ge3$ is an odd integer, then
\[
 \mathcal B(n)=\pi,
 \qquad \mu_n=\tfrac12\delta_{-\pi/2}+\tfrac12\delta_{\pi/2}
\]
is the unique minimizer up to translation.
\item Otherwise write
\[
 n=2k+1<r<n+2,\qquad k\in\mathbb Z_{\ge0},\qquad s=(-1)^k.
\]
When $k=0$, put $t=0$ and $M=1$.  When $k\ge1$, let $t$ be the unique
solution in
\[
 \frac{k\pi}{r}<t<\frac{(k+1/2)\pi}{r}
\]
of
\begin{equation}\label{filt:tangent}
 \tan t=r\tan(rt-k\pi),
 \qquad M=s\frac{\cos(rt)}{\cos t}>0.
\end{equation}
There is a unique $T$ in
\[
 \frac\pi2<T<\frac{(n+2)\pi}{2r}
\]
satisfying
\begin{equation}\label{filt:crossing}
 s\frac{\cos(rT)}{\cos T}=M.
\end{equation}
Then $\mathcal B(r)=2T$.  Put
\[
 P=\frac{-\cos T}{\cos t-\cos T},\qquad
 Q=\frac{\cos t}{\cos t-\cos T}.
\]
The unique minimizer up to translation is
\begin{equation}\label{filt:measure}
 \mu_r=\frac P2(\delta_{-t}+\delta_t)
             +\frac Q2(\delta_{-T}+\delta_T).
\end{equation}
At $t=0$ the two coincident atoms in the first term are combined.
\item In the first band $1<r<3$, these formulas simplify to
\begin{equation}\label{filt:firstband}
 T=\frac{2\pi}{1+r},\qquad
 \mathcal B(r)=\frac{4\pi}{1+r},\qquad
 P=\frac{-c}{1-c},\quad Q=\frac1{1-c},\quad
 c=\cos\frac{2\pi}{1+r}.
\end{equation}
\end{enumerate}
For arbitrary frequencies $0<\alpha<\gamma$, the least diameter is
$\alpha^{-1}\mathcal B(\gamma/\alpha)$, and the atoms are obtained by
dividing the coordinates above by $\alpha$.
\end{theorem}

\begin{theorem}[Asymmetric resonance law]\label{filt:cusp}
Fix an odd integer $n\ge3$, and put
\[
 A_n=\frac{n(n^2-1)}6,
 \qquad
 C_n=\left(\frac{3\pi}{2n(n^2-1)}\right)^{1/3}.
\]
As $\varepsilon\downarrow0$, the optimal inner and outer coordinates for
$r=n+\varepsilon$ satisfy
\begin{align}
 t&=\frac\pi2-C_n\varepsilon^{1/3}+O_n(\varepsilon),\label{filt:innercusp}\\
 T&=\frac\pi2+2C_n\varepsilon^{1/3}+O_n(\varepsilon),\label{filt:outercusp}\\
 \mathcal B(n+\varepsilon)
   &=\pi+4C_n\varepsilon^{1/3}+O_n(\varepsilon).\label{filt:widthcusp}
\end{align}
Furthermore $P=2/3+O_n(\varepsilon^{2/3})$ and
$Q=1/3+O_n(\varepsilon^{2/3})$.

For the approach from below, let $s_n=(-1)^{(n-1)/2}$ and
\[
 M_n^- =\max_{0\le u<\pi/2}
                 \left(-s_n\frac{\cos(nu)}{\cos u}\right).
\]
The maximum exists and is finite.  Then
\begin{equation}\label{filt:left}
 \mathcal B(n-\varepsilon)
 =\pi+\frac{\pi}{n+M_n^-}\varepsilon+O_n(\varepsilon^2).
\end{equation}
In particular $M_3^-=1$ and
$\mathcal B(3-\varepsilon)=\pi+\pi\varepsilon/4+O(\varepsilon^2)$,
whereas the right-hand increase has order $\varepsilon^{1/3}$.
\end{theorem}

The cube-root law concerns perturbation of the prescribed frequencies, not
the location of zeros of any particular $L$-function.  Its assertion is a
sharp design constraint: moving the higher frequency just above an odd
multiple of the lower one requires a much larger increase in the averaging
window than a perturbation of the same size in the opposite direction.

\section{Proof of the classification}

\subsection{A scalar quotient and its last positive lobe}

Fix a nonresonant ratio $r$ and $n,k,s$ as in
Theorem~\ref{filt:classification}.  On $[0,\pi/2)$ define
\[
 f(u)=s\frac{\cos(ru)}{\cos u}.
\]
The maximum of $f$ is positive, finite, and attained.  Indeed,
$s\cos(r\pi/2)<0$, so $f(u)\to-\infty$ at the right endpoint.

When $k=0$, the only positive lobe before $\pi/2$ is
$0\le u<\pi/(2r)$, and $\cos(ru)\le\cos u$ there.  Consequently
$\max f=f(0)=1$, with equality only at $u=0$.

Suppose $k\ge1$.  The last positive lobe is
\[
 J=\left(\frac{(k-1/2)\pi}{r},
               \frac{(k+1/2)\pi}{r}\right).
\]
Every point of an earlier positive lobe can be shifted forward by an integer
multiple of $2\pi/r$ into $J$.  This leaves its positive numerator unchanged
and strictly decreases its positive denominator.  Thus the global maximum
occurs in $J$.  Before $k\pi/r$ the function is strictly increasing.

For completeness, let
\[
 W(u)=s\{-r\sin(ru)\cos u+\cos(ru)\sin u\}.
\]
Then $f'(u)=W(u)/\cos^2u$ and
\begin{equation}\label{filt:Wprime}
 W'(u)=(1-r^2)s\cos(ru)\cos u.
\end{equation}
It is strictly negative throughout $J$.  At $u=k\pi/r$ the value of $W$ is
positive; at the right endpoint of $J$ it is negative.  There is exactly one
stationary point $t$, and it is the unique global maximum.  Dividing
$W(t)=0$ by the nonzero cosines gives~\eqref{filt:tangent}.
We have proved
\begin{equation}\label{filt:max}
 M=\max_{0\le u<\pi/2}f(u),
 \qquad \{u:f(u)=M\}=\{t\}.
\end{equation}

Now continue $f$ to
\[
 \pi/2<u<U:=\frac{(n+2)\pi}{2r}.
\]
Here $s\cos(ru)<0$ and $\cos u<0$, so $f>0$; its endpoint limits are
$+\infty$ and $0$.  Equation~\eqref{filt:Wprime} again gives $W'<0$ on this
interval.  Its value just after $\pi/2$ is
$s\cos(r\pi/2)<0$, hence $f'<0$ throughout.  This proves the existence and
uniqueness of $T$ in~\eqref{filt:crossing}.

\subsection{A nonnegative certificate and the endpoint measure}

Define the even function
\begin{equation}\label{filt:g}
 g(u)=M\cos u-s\cos(ru).
\end{equation}
For $0\le u<\pi/2$, equations~\eqref{filt:max} give $g(u)\ge0$.
For $\pi/2<u<T$ we have $f(u)>M$ and $\cos u<0$, whence again
$g(u)=\cos u\{M-f(u)\}>0$.
At $\pi/2$ the value is $-s\cos(r\pi/2)>0$.
It follows that
\begin{equation}\label{filt:certzeros}
 g\ge0\text{ on }[-T,T],\qquad
 \{g=0\}=\{-T,-t,t,T\},
\end{equation}
with the two inner entries identified when $t=0$.
The inner zeros are double and the endpoint zeros simple:
\[
 g''(t)=(r^2-1)M\cos t>0,
 \qquad g'(T)=-\cos T\,f'(T)<0.
\]

Let an admissible measure have support diameter at most $2T$ and translate
it into $[-T,T]$.  The vanishing Fourier moments imply
$\int g\,d\mu=0$.  By~\eqref{filt:certzeros}, its support lies in the listed
finite set.  The measure~\eqref{filt:measure} is a probability measure with
positive weights and zero cosine moments, because
\[
 P\cos t+Q\cos T=0,
 \qquad \cos(rt)=sM\cos t,\quad
 \cos(rT)=sM\cos T.
\]
Symmetry gives its zero sine moments.

We next prove that there are no nonsymmetric minimizers.  For $t>0$ let
$d_t$ and $d_T$ be the differences of the masses at the positive and negative
inner and outer nodes.  The imaginary parts of the two vanishing equations
read
\[
 \begin{pmatrix}\sin t&\sin T\\\sin(rt)&\sin(rT)\end{pmatrix}
 \begin{pmatrix}d_t\\d_T\end{pmatrix}=0.
\]
At the inner node, $r\sin(rt)=sM\sin t$.  At the outer node,
$W(T)<0$ and $\cos T<0$, which give
\[
 M\sin T-sr\sin(rT)>0.
\]
Thus
\[
 \sin t\,\sin(rT)-\sin T\,\sin(rt)
 =\sin t\left\{\sin(rT)-\frac{sM}{r}\sin T\right\}\ne0.
\]
Both differences are zero.  For $t=0$, the first sine equation already
forces equal endpoint masses.  Total mass and the first cosine equation
then uniquely determine $P,Q$.  The endpoint masses are positive, so no
admissible measure fits in a strictly shorter interval.  This proves the
nonresonant cases.

For odd $r=n$, the function $\cos u$ is nonnegative on
$[-\pi/2,\pi/2]$.  Its integral can vanish only for a measure supported at
the two endpoints.  The sine moment forces equal masses, and that measure
also annihilates frequency $n$.  Any smaller interval has strictly positive
cosine after centering.  This proves the resonant case.

Finally, when $1<r<3$,
\[
 \cos u-\cos(ru)
 =2\sin\frac{(1+r)u}{2}\sin\frac{(r-1)u}{2}.
\]
The first positive zero beyond the origin is $T=2\pi/(1+r)$; here
$\cos T=\cos(rT)<0$.  This yields~\eqref{filt:firstband} and completes the
proof.\hfill$\square$

\section{Proof of the asymmetric resonance law}

\subsection{Localization at an odd resonance}

Fix odd $n\ge3$, put $r=n+\varepsilon$ and
$h=\pi\varepsilon/2$.  Write the optimal inner node as
$t=\pi/2-y$.  The quotient at that node is
\begin{equation}\label{filt:innerquotient}
 M=\frac{\sin(ry-h)}{\sin y}.
\end{equation}
At $\varepsilon=0$, the limiting quotient is $\sin(ny)/\sin y$.
For $0<y\le\pi/2$ it is strictly less than $n$.  One way to see the strict
bound is to write $\sin(ny)/\sin y$ as a sum of $n$ complex numbers of
modulus one; equality in the triangle inequality would require all adjacent
phases $e^{2iy}$ to coincide, which cannot happen in this interval.
Testing~\eqref{filt:innerquotient} at $y=\varepsilon^{1/3}$ shows that its
maximum tends to $n$.  Uniform convergence on compact sets away from $y=0$
therefore forces $y\to0$.

The stationary equation is
\begin{equation}\label{filt:stationarynear}
 H(\varepsilon,y):=
 r\cos(ry-h)\sin y-\sin(ry-h)\cos y=0.
\end{equation}
As $(\varepsilon,y)\to(0,0)$ with $\varepsilon,y\ge0$,
\[
 H(\varepsilon,y)
 =\frac\pi2\varepsilon-2A_ny^3
 +O_n\bigl(y^5+\varepsilon y^2+\varepsilon y^3
                 +\varepsilon^2y+\varepsilon^3\bigr).
\]
The positive solution thus has $y\asymp_n\varepsilon^{1/3}$.  More explicitly,
the remainder is $o(\varepsilon+y^3)$ at the solution, so
$y^3/\varepsilon\to\pi/(4A_n)=C_n^3$.

The stated sharper remainder follows from an elementary analytic implicit
function argument.  Put $z=\varepsilon^{1/3}$ and $y=zv$.  The function
$H(z^3,zv)/z^3$ extends real analytically to $z=0$ and equals
\[
 \frac\pi2-2A_nv^3+O_n(z^2)
\]
for $v$ near $C_n$.  Its derivative with respect to $v$ at $(0,C_n)$ is
$-6A_nC_n^2\ne0$.  Hence $v=C_n+O_n(z^2)$, and
$y=C_n\varepsilon^{1/3}+O_n(\varepsilon)$.

\subsection{The outer node and the factor two}

Write $T=\pi/2+x$.  The crossing equation becomes
\begin{equation}\label{filt:nearcompare}
 \frac{\sin(rx+h)}{\sin x}
 =\frac{\sin(ry-h)}{\sin y}.
\end{equation}
For $r=n$ the left quotient is less than $n$ whenever $x>0$ is fixed,
while $M\to n$.  Monotonicity of the crossing branch therefore gives
$x\to0$.

Set $x=zw$ and use the just-constructed analytic solution $v(z)$ for
$y=zv(z)$.  After subtracting the two sides
of~\eqref{filt:nearcompare} and dividing by $z^2$, the equation extends
analytically to $z=0$ as
\begin{equation}\label{filt:reducedcubic}
 A_n\left(\frac{2C_n^3}{w}-w^2+3C_n^2\right)=0.
\end{equation}
Its positive solution is $w=2C_n$, because multiplication by $w$ factors
the corresponding cubic as
$(w-2C_n)(w+C_n)^2=0$.
The derivative at $2C_n$ is $-9A_nC_n/2\ne0$.
The analytic equation before taking $z=0$ differs from the displayed
one by $O_n(z^2)$, uniformly for $w$ near $2C_n$. Indeed,
\[
 \frac{\sin((n+z^3)zw\pm\pi z^3/2)}{\sin(zw)}
 =n+z^3+z^2\left(\pm\frac{\pi}{2w}-A_nw^2\right)+O_n(z^4),
\]
and the common $z^3$ term cancels when the two quotients are subtracted;
also $v(z)=C_n+O_n(z^2)$.
The implicit function theorem therefore gives a positive solution
$w=2C_n+O_n(z^2)$.  The uniqueness already proved for the crossing equation
identifies it with the optimal $T$.  This proves
\eqref{filt:innercusp}--\eqref{filt:widthcusp}.
The weight formulas and
$\cos t=\sin y$, $-\cos T=\sin x$ give
$P=2/3+O_n(z^2)$ and $Q=1/3+O_n(z^2)$.

\subsection{The regular approach from below}

For $r=n-\varepsilon$, the sign of the preceding band is $-s_n$.
The limiting maximum $M_n^-$ occurs strictly before $\pi/2$, since the
quotient tends to $-n$ there.  For $n=3$ the maximum is at $0$ and equals
$1$.  For $n\ge5$ the preceding last-lobe argument gives a unique interior
maximum with strictly negative second derivative.  In either case its
maximum is $M_n^-+O_n(\varepsilon)$.
With $T=\pi/2+x$ and $h=\pi\varepsilon/2$, the outer equation now reads
\[
 -\frac{\sin((n-\varepsilon)x-h)}{\sin x}
      =M_n^-+O_n(\varepsilon).
\]
Setting $x=\varepsilon v$ extends the equation analytically to
$\pi/(2v)-n=M_n^-$.  The nonzero derivative in $v$ yields
\[
 x=\frac{\pi}{2(n+M_n^-)}\varepsilon+O_n(\varepsilon^2),
\]
which is~\eqref{filt:left}.\hfill$\square$

\section{Robustness and the exact first-order cost of a shorter window}

For fixed nonresonant $r$ write $t,T,M,P,Q,g$ as above and let
\[
 \epsilon(\nu)=\max\{|\widehat\nu(1)|,|\widehat\nu(r)|\}.
\]
Here $W_1$ is the usual transport distance with cost $|u-v|$.

\begin{proposition}[Sharp stability exponent]\label{filt:stability}
There is $C_r<\infty$ such that every probability measure $\nu$ on
$[-T,T]$ satisfies
\[
 W_1(\nu,\mu_r)\le C_r\sqrt{\epsilon(\nu)}.
\]
No exponent greater than $1/2$ can hold uniformly for that fixed $r$.
At an odd resonance, on the optimal interval $[-\pi/2,\pi/2]$, the optimal
exponent is $1$ instead.
\end{proposition}

\begin{proof}
Let $G$ be the finite zero set in~\eqref{filt:certzeros}.  The double inner
zeros and simple endpoint zeros imply
$g(u)\ge c_r\operatorname{dist}(u,G)^2$ on the interval.  Also
$\int g\,d\nu\le(M+1)\epsilon(\nu)$.
Round every point to a nearest node of $G$.  Cauchy--Schwarz bounds the
transport cost by $O_r(\sqrt{\epsilon})$, and the Fourier moments change
by the same order.  Total mass and the two complex moments determine the
weights on $G$ uniquely: this is exactly the even and odd linear algebra
in the uniqueness proof.  Their inverse on its image is bounded, giving
the asserted estimate.  Splitting an inner atom into equal masses at
distance $a$ on either side changes both moments by $O_r(a^2)$ but has
transport distance bounded below by a positive multiple of $a$.  Thus the
exponent cannot be improved.

At an odd resonance, $\cos u$ is bounded below by a positive constant times
the distance to the endpoints.  Rounding therefore costs $O(\epsilon)$,
and the first sine moment controls the imbalance of the endpoint weights.
Moving the endpoints inward by $a$ proves sharpness.
\end{proof}

\begin{theorem}[First-order shortening law]\label{filt:shortening}
For $d>0$ sufficiently small take the infimum below over probability
measures $\nu$, and put
\[
 D_r(d)=\inf_{\operatorname{supp}\nu\subset[-T+d/2,T-d/2]}
             \epsilon(\nu),
 \qquad
 L_r=\frac{Q[-g'(T)]}{2(M+1)}>0.
\]
Then
\begin{equation}\label{filt:shorteninglaw}
 D_r(d)=L_rd+O_r(d^2).
\end{equation}
Moreover at $r=n+\varepsilon$, $n\ge3$ odd,
\begin{equation}\label{filt:leakcusp}
 L_{n+\varepsilon}
  =\frac{3A_nC_n^2}{2(n+1)}\varepsilon^{2/3}
          +O_n(\varepsilon).
\end{equation}
\end{theorem}

\begin{proof}
Put $a=-g'(T)/2$, $K=Qa$, and
$h_0(u)=K\cos u/\cos t$.  At the inner nodes $h_0=K$, while at the
outer nodes $h_0=K-a$.  On the shortened interval,
\[
 \min\{g(u)+dh_0(u)\}\ge dK-O_r(d^2).
\]
Indeed, near an inner zero the positive quadratic term can lose only
$O(d^2)$ under a linear perturbation, and at each moved endpoint the
linear inward increase is $da$; away from these finitely many neighborhoods
the unperturbed function is bounded below by a positive constant.
Integrating and using the definition of $\epsilon$ gives
\[
 (M+1+dK/\cos t)\epsilon(\nu)\ge dK-O_r(d^2),
\]
the required lower estimate.

For the matching upper bound, keep the inner nodes at $\pm t$, move the
outer nodes to $\pm(T-d/2)$, and choose their two symmetric total weights
so that
$\widehat\nu(1)+s\widehat\nu(r)=0$.
Such positive weights exist for small $d$, since at $d=0$ the two values
of $\cos u+s\cos(ru)$ have opposite signs.  Both moments are real and
$\widehat\nu(r)=-s\widehat\nu(1)$; therefore
\[
 (M+1)\widehat\nu(1)=\int g\,d\nu
                  =dQa+O_r(d^2).
\]
This proves~\eqref{filt:shorteninglaw}.

For~\eqref{filt:leakcusp}, write $y=C_n\varepsilon^{1/3}+O_n(\varepsilon)$
and $x=2y+O_n(\varepsilon)$.  The expansion near the outer endpoint is
\[
 g(\pi/2+x)=h+3A_ny^2x-A_nx^3+O_n(\varepsilon^{4/3}),
 \qquad h=2A_ny^3+O_n(\varepsilon^{5/3}).
\]
Differentiating the underlying analytic trigonometric expression gives
$-g'(T)=9A_nC_n^2\varepsilon^{2/3}+O_n(\varepsilon)$.
Now use $Q=1/3+O_n(\varepsilon^{2/3})$ and $M=n+O_n(\varepsilon^{2/3})$.
\end{proof}

\section{An exact nonconsecutive example}

The ratio $r=4$ permits a radical expression.  Put
\[
 x=\sqrt{\frac{2+\sqrt{10}}{12}},\qquad
 z=-x+\sqrt{1-2x^2}.
\]
Then
\[
 t=\arccos x,\qquad T=\arccos z,\qquad
 \boxed{\mathcal B(4)=2\arccos\left(
 -\sqrt{\frac{2+\sqrt{10}}{12}}
 +\sqrt{\frac{4-\sqrt{10}}6}\right)}.
\]
To verify this without numerical root finding, use
$\cos(4u)=8\cos^4u-8\cos^2u+1$.  Tangency of
$-\cos(4u)/\cos u$ gives $24x^4-8x^2-1=0$, and, with
$M=-(8x^4-8x^2+1)/x$, the polynomial identity
\[
 8v^4-8v^2+1+Mv
   =8(v-x)^2(v^2+2xv+3x^2-1)
\]
identifies $z$ as the outer crossing.  The other quadratic root is less
than $-1$, whereas $-1<z<0$.  Numerically,
\[
 \begin{array}{c|r}
 t&0.8554371830421481\\
 T&1.8569146032739532\\
 \mathcal B(4)&3.7138292065479064\\
 P&0.3008471645880608\\
 Q&0.6991528354119392
 \end{array}
\]
The exact radical formulas, not these displayed decimal approximations,
define the example.

The constant on the other side of the next resonance is also elementary:
\[
 M_5^-=-\min_{0\le u<\pi/2}
         \frac{\cos(5u)}{\cos u}
       =-\min_{0\le v\le1}(16v^4-20v^2+5)=\frac54.
\]
Consequently
$\mathcal B(5-\varepsilon)=\pi+(4\pi/25)\varepsilon
+O(\varepsilon^2)$, alongside the cube-root increase to the right of $5$.

\section{Positive dilation averages for zeros on a common vertical line}

The width theorem is unchanged if the two constraints have a common real
part.  More precisely, let $\beta\in\mathbb R$ and $0<\alpha<\gamma$.
Among probability measures satisfying
\[
 \int e^{(\beta+i\alpha)u}\,d\mu(u)=
 \int e^{(\beta+i\gamma)u}\,d\mu(u)=0,
\]
the least support diameter is $\alpha^{-1}\mathcal B(\gamma/\alpha)$.
Indeed, exponential tilting
\[
 d\nu(u)=\frac{e^{\beta u}\,d\mu(u)}{\int e^{\beta v}\,d\mu(v)}
\]
preserves support and transforms the conditions into ordinary Fourier
vanishing.  Conversely, tilt the optimal measure back by $e^{-\beta u}$.
If its unweighted nodes and masses are $u_j,w_j$, the required positive
masses are
\begin{equation}\label{filt:tiltweights}
 \widetilde w_j=
       \frac{w_je^{-\beta u_j}}{\sum_\ell w_\ell e^{-\beta u_\ell}}.
\end{equation}

Under the Dirichlet GRH theorem supplied by the upstream Lean development,
two distinct positive ordinates $\alpha,\gamma$ of the same or different
Dirichlet $L$-functions lie on the common line $\beta=1/2$.  In any explicit
formula in which a zero contributes $c_\rho x^\rho$, averaging the quantity
at $xe^{u_j}$ with the positive weights~\eqref{filt:tiltweights} multiplies
that contribution by
$\sum_j\widetilde w_je^{\rho u_j}$.  Our optimal average cancels the selected
frequencies exactly.  For example, with a smooth compactly supported
function $W$ on $(0,\infty)$, each zero term in the smoothed von Mangoldt
sum $\sum_m\Lambda(m)\chi(m)W(m/x)$ has this form.  The operation remains a
positive average on the original sum and uses the least possible
multiplicative scale ratio
\[
 \exp\!\left(\frac{\mathcal B(\gamma/\alpha)}{\alpha}\right).
\]
Multiplicity of a selected zero changes its coefficient and does not
affect this simple cancellation.  Requiring vanishing derivatives of the
multiplier is a different, stronger design problem.

This theorem resolves the two-frequency case of the merged article's
question about nonconsecutive prescribed frequencies.  It supplements the
article's exact equally spaced and consecutive-harmonic formulas.  The
optimization theorem and the resonance law are unconditional statements
about positive measures; Dirichlet GRH is used only to place the selected
arithmetic zeros on one common vertical line.

\section{Relation to prior work and novelty limits}

The mathematical setting belongs to the classical extremal theory of
positive-definite entire functions.  In particular, Logan studied the
maximum number of zeros in an interval for positive-definite bandlimited
functions [B.~F.~Logan, \emph{SIAM Journal on Mathematical Analysis}
\textbf{14} (1983), 258--268, DOI:
\href{https://doi.org/10.1137/0514024}{10.1137/0514024}].
The same optimization has a long engineering history as positive
two-mode input shaping: an impulse sequence cancels two undamped vibration
modes.  Crain, Singhose and Seering studied simultaneous two-mode shapers
in their 1996 IFAC paper, and the three-impulse formula in the first band
is already part of that literature.  A primary historical account is
\href{https://web.mit.edu/shaping/www/papers/2MIFAC2173.ps}{their author-hosted
paper}; related explicit formulas also appear in
\href{https://patents.google.com/patent/US5638267A/en}{US Patent 5,638,267}.
No novelty is claimed for the idea of positive two-frequency cancellation
or for~\eqref{filt:firstband}.

The merged article already proves optimal positive filters for a single
frequency with jets and for blocks of consecutive harmonics.  The present
note supplies a self-contained global proof over all positive probability
measures, proves uniqueness without assuming symmetry or a prescribed
number of atoms, and derives the asymmetric cubic-root law at every odd
resonance together with its exact quantitative shortening consequence.
The resonance and shortening laws are the principal candidate original
contributions; the global classification supplies their rigorous basis.

A targeted search of primary-source records did not locate the cubic-root
law or the stated shortening asymptotic.  The engineering literature
contains substantial prior work on the optimizing curves, so priority
for the global classification is not asserted.  This is a limited novelty
check, not proof of historical priority.  All the mathematical claims
above have self-contained proofs, independent of any novelty claim.
